% Part: lambda-calculus
% Chapter: introduction
% Section: minimization

\documentclass[../../../include/open-logic-section]{subfiles}

\begin{document}

\olfileid[tr]{lam}{int}{min}
\olsection{\usetoken{S}{lambda definable} Fonksiyonlar Minimizasyon Altında Kapalıdır}

\begin{lem}
$f(x,y)$'nin !!{lambda definable} olduğunu varsayın. $g$,
\[
g(x) \simeq \umin{y}{f(x,y)}
\]
ile tanımlansın. O zaman $g$ !!{lambda definable}.
\end{lem}

\begin{proof}
Düşünce kabaca şöyledir. $x$ verildiğinde, $n$'den başlayarak bir $y$ arayan
$h_x(n)$ fonksiyonunu tanımlamak için sabit nokta lambda terimi $Y$'yi
kullanacağız; o zaman $g(x)$ yalnızca $h_x(0)$'dır. $h_x$ fonksiyonu bir
sabit nokta denkleminin çözümü olarak ifade edilebilir:
\[
h_x(n) \simeq
\begin{cases}
n & \text{$f(x,n) = 0$ ise} \\
h_x(n+1) & \text{değilse.}
\end{cases}
\]

Ayrıntılar şöyledir. $f$ ilkel özyinelemeli olduğundan bir $F$ terimi
tarafından !!{lambda defined}. Ayrıca $D(M,N,\bar{0})\red M$ ve
$D(M,N,\bar{1})\red N$ olacak biçimde bir $D$ lambda terimimiz olduğunu
anımsayın. Şimdilik $x$'i sabitleyerek $h_x$'i !!{lambda define} etmek için
$x$'e bağlı, şu koşulu taşıyan bir $H$ terimi bulmak isteriz:
\[
H(\num{n}) \equiv D(\num{n}, H(S(\num{n})), F(x, \num{n})).
\]
Bunu sabit nokta terimi $Y$ ile yapabiliriz. Önce $U$,
\[
\lambd[h][\lambd[z][D(z,(h(Sz)),F(x,z))]]
\]
terimi olsun, ardından $H$, $YU$ terimi olsun. $H$ içindeki tek serbest
değişkenin $x$ olduğuna dikkat edin. $H$'nin yukarıdaki denklemi sağladığını
gösterelim.

$Y$'nin tanımı gereği
\[
H = YU \equiv U(YU) = U(H).
\]
Özellikle her $n$ doğal sayısı için
\begin{align*}
H(\num{n}) & \equiv U(H, \num{n}) \\
& \red D(\num{n}, H(S(\num{n})), F(x, \num{n})),
\end{align*}
istendiği gibi. Son satırda $x$ yerine $\num{m}$ koyarsanız, ifade
$F(\num{m},\num{n})$, $\num{0}$'a indirgeniyorsa $\num{n}$'ye; başka
herhangi bir sayıya indirgeniyorsa $H(S(\num{n}))$'ye indirgenir.

İspatı bitirmek için $G$, $\lambd[x][H(\num 0)]$ olsun. O zaman $G$, $g$'yi
!!{lambda define}s; başka bir deyişle her $m$ için $G(\num m)$, $g(m)$
tanımlıysa $\overline{g(m)}$'ye indirgenir, değilse normal biçimi yoktur.
\end{proof}

\end{document}
